\begin{abstract}
The agent harness---the system layer comprising prompts, tools, memory, and orchestration logic that surrounds the model---has emerged as the central engineering abstraction for LLM-based agents. Yet harness design remains ad~hoc, with no formal theory governing composition, preservation of properties under compilation, or systematic comparison across frameworks. We show that the categorical Architecture triple $(G, \mathrm{Know}, \Phi)$ from the ArchAgents framework provides exactly this formalization. The four pillars of agent externalization (Memory, Skills, Protocols, Harness Engineering) map onto the triple's components: Memory as coalgebraic state, Skills as operad-composed objects, Protocols as syntactic wiring~$G$, and the full Harness as the Architecture itself. Structural guarantees---integrity gates, quality-based escalation, convergence detection---are $\mathrm{Know}$-level properties that transfer under compilation via functor laws (Proposition~5.1). We validate this correspondence with a reference implementation featuring compiler functors targeting Swarms, DeerFlow, Ralph, and Scion (with full architectural preservation) and LangGraph (certificate-only preservation via a single-node wrapper), achieving 100\% certificate preservation across all targets. An end-to-end escalation experiment with real LLM agents confirms that quality-based model escalation---a harness-level guarantee---operates independently of which model executes. The result positions categorical architecture as the formal theory behind harness engineering.
\end{abstract}
